\documentclass[11pt,a4paper]{article}

\usepackage[T1]{fontenc}
\usepackage[utf8]{inputenc}
\usepackage[french]{babel}
\usepackage{lmodern}
\usepackage[a4paper,margin=2.4cm]{geometry}
\usepackage{graphicx}
\usepackage{xcolor}
\usepackage{titlesec}
\usepackage{enumitem}
\usepackage{tcolorbox}
\usepackage{fancyhdr}
\usepackage{array}
\usepackage{longtable}
\usepackage{booktabs}
\usepackage[hidelinks]{hyperref}

% Palette Travel's Affair
\definecolor{taSky}{HTML}{00ADEF}
\definecolor{taNavy}{HTML}{1B2A7A}
\definecolor{taDark}{HTML}{0E1A4A}
\definecolor{taRed}{HTML}{E4002B}
\definecolor{taMist}{HTML}{EEF3FB}
\definecolor{taLine}{HTML}{D0D8E8}
\definecolor{taStone}{HTML}{5B6472}

\hypersetup{colorlinks=true,linkcolor=taNavy,urlcolor=taSky}

% Titres
\titleformat{\part}{\color{taRed}\huge\bfseries}{Partie \thepart}{0.8em}{}
\titleformat{\section}{\color{taNavy}\Large\bfseries}{\thesection.}{0.6em}{}
\titleformat{\subsection}{\color{taDark}\large\bfseries}{\thesubsection}{0.6em}{}
\titlespacing*{\section}{0pt}{1.4em}{0.6em}

% En-tete / pied
\pagestyle{fancy}
\fancyhf{}
\renewcommand{\headrulewidth}{0.4pt}
\renewcommand{\footrulewidth}{0.4pt}
\fancyhead[L]{\small\color{taNavy}\textbf{Travel's Affair}}
\fancyhead[R]{\small\color{taStone}Guide de la plateforme}
\fancyfoot[L]{\small\color{taStone}Documentation fonctionnelle}
\fancyfoot[R]{\small\color{taStone}\thepage}

% Encadres
\newtcolorbox{astuce}{colback=taMist,colframe=taSky,boxrule=0.8pt,arc=3pt,left=8pt,right=8pt,top=6pt,bottom=6pt}
\newtcolorbox{attention}{colback=taRed!6,colframe=taRed,boxrule=0.8pt,arc=3pt,left=8pt,right=8pt,top=6pt,bottom=6pt}

\newcommand{\etat}[1]{\textcolor{taNavy}{\texttt{#1}}}
\newcommand{\chemin}[1]{\textcolor{taStone}{\texttt{#1}}}

\begin{document}

% ====================== COUVERTURE ======================
\begin{titlepage}
\centering
\vspace*{1.6cm}
\IfFileExists{../public/img/site/logo.jpg}{\includegraphics[width=5cm]{../public/img/site/logo.jpg}\\[1.0cm]}{}
{\color{taNavy}\Huge\bfseries Guide de la plateforme\par}
\vspace{0.4cm}
{\color{taSky}\Large Espace client \& portail d'administration\par}
\vspace{1.8cm}
\begin{tcolorbox}[colback=taMist,colframe=taNavy,width=0.84\textwidth,arc=4pt]
\centering\color{taDark}
\large Plateforme de réservation \textbf{Travel's Affair}\\[4pt]
\normalsize Vols \textbullet{} Hébergements \textbullet{} Transports \textbullet{} Location de véhicules \textbullet{} Activités\\[4pt]
\small Voyages Afrique $\leftrightarrow$ Europe
\end{tcolorbox}
\vfill
{\color{taStone}\small Agence de Yaoundé, Face Tsinga Dubaï, Immeuble Italia Optique\\
875 Av. du 27 Août 1940, BP 13368, Yaoundé, Cameroun\\[4pt]
+237 696 35 82 62 \textbullet{} agence.yaounde@travels-affair.com\par}
\vspace{0.5cm}
{\color{taStone}\small \today}
\end{titlepage}

\tableofcontents
\newpage

% ##########################################################
\part{Présentation générale}
% ##########################################################

\section{La plateforme en bref}
Travel's Affair est une plateforme de réservation de voyages couvrant l'axe
\textbf{Afrique $\leftrightarrow$ Europe}. Elle réunit cinq familles de services, un parcours
de réservation et de paiement en ligne, un espace client, et un portail d'administration
offrant une \textbf{vision 360\textdegree} de toute l'activité.

\subsection{Les cinq services}
\begin{center}
\begin{longtable}{>{\raggedright}p{3.4cm} p{10cm}}
\toprule
\textbf{Service} & \textbf{Description} \\
\midrule
\endhead
\textbf{Vols} & Aller simple, aller-retour ou multi-destinations ; escales, franchise bagages, places disponibles. \\
\textbf{Hébergements} & Hôtels, résidences, lodges, appartements, maisons d'hôtes ; réservation par bien (chambre/suite) et par dates. \\
\textbf{Transports} & Trajets terrestres assurés par des prestataires référencés. \\
\textbf{Location de véhicules} & Véhicules par marque, modèle et catégorie, sur une période donnée. \\
\textbf{Activités} & Excursions et activités à date programmée. \\
\bottomrule
\end{longtable}
\end{center}

\subsection{Langues et devises}
\begin{itemize}[leftmargin=1.4em]
  \item \textbf{Langues :} français (par défaut) et anglais. Le choix est mémorisé pour le visiteur.
  \item \textbf{Devises :} \etat{XAF} (par défaut), \etat{EUR}, \etat{USD}.
\end{itemize}

\begin{astuce}
La langue par défaut est le \textbf{français} et la devise par défaut le \textbf{franc CFA
(XAF)} : indépendamment des réglages du navigateur. Le visiteur peut en changer à tout moment.
\end{astuce}

\subsection{Identité visuelle}
La plateforme reprend la charte de l'agence : bleu ciel \textbf{\#00ADEF}, bleu marine
\textbf{\#1B2A7A} et rouge d'action \textbf{\#E4002B}, logo et photos de l'agence, logos des
compagnies partenaires. Un widget de \textbf{chat en ligne} (Tawk.to) est disponible sur le site.

% ##########################################################
\part{Espace client (site public)}
% ##########################################################

\section{Navigation du site}
Le site public donne accès aux rubriques suivantes :
\begin{itemize}[leftmargin=1.4em]
  \item \textbf{Accueil} : recherche, « Comment ça marche », services, offres en vedette, vitrine de l'agence et partenaires.
  \item \textbf{Réserver} (menu déroulant) : Vols, Hébergements, Transports, Location, Activités.
  \item \textbf{L'agence}, \textbf{Infos pratiques}, \textbf{FAQ}, \textbf{Contact}.
  \item \textbf{Compte} : connexion, inscription, mot de passe oublié, mon compte.
\end{itemize}

\section{Parcours de réservation}
\begin{enumerate}[leftmargin=1.6em]
  \item \textbf{Rechercher / parcourir} un service et ouvrir une offre pour en voir le détail
        (prix, conditions, politique d'annulation).
  \item \textbf{Réserver} : renseigner les voyageurs (et les dates/quantités pour un hébergement
        ou une location). Un compte est requis ; sinon le visiteur est invité à se connecter.
  \item \textbf{Payer} : la réservation est créée en \etat{PENDING\_PAYMENT} et redirige vers
        l'écran de paiement.
\end{enumerate}

\section{Moyens de paiement}
Selon la configuration de l'agence, les moyens suivants sont pris en charge :
\begin{center}
\begin{longtable}{>{\raggedright}p{4.6cm} p{8.8cm}}
\toprule
\textbf{Moyen} & \textbf{Prestataire} \\
\midrule
\endhead
Carte bancaire & Stripe \\
PayPal & PayPal \\
Orange Money / MTN MoMo & Campay (ou CinetPay) \\
Express Union & CinetPay \\
(Mode démonstration) & Paiement simulé si aucun prestataire n'est configuré \\
\bottomrule
\end{longtable}
\end{center}
Après paiement, le client est redirigé vers la plateforme ; les prestataires notifient le
serveur (webhooks) pour confirmer l'encaissement et faire évoluer la réservation.

\section{Espace « Mon compte »}
Le client y retrouve ses réservations, leur statut, et le détail de chaque dossier
(\chemin{/mon-compte/reservations/[id]}) : récapitulatif, paiements, documents et possibilité
d'\textbf{annuler} une réservation encore en attente de paiement.

\section{Vérification d'un dossier (QR code)}
Chaque billet/bon porte un \textbf{QR code} renvoyant vers une page publique de vérification
(\chemin{/verifier/[référence]}). Elle affiche un état :
\begin{itemize}[leftmargin=1.4em]
  \item \textbf{Valide} : statut payé, en attente de confirmation, confirmé ou modifié ;
  \item \textbf{Non réglé} : en attente de paiement ;
  \item \textbf{Invalide} : annulé/remboursé ou référence inconnue.
\end{itemize}

\begin{attention}
\textbf{Ce n'est pas un billet de voyage.} Les documents émis par la plateforme (confirmation,
récapitulatif, bon de réservation) attestent la réservation mais ne constituent pas le titre de
transport. Le \textbf{billet (ou bon) de voyage définitif est délivré par l'agence} Travel's
Affair après confirmation. Cette mention est rappelée sur les PDF, dans l'espace client et sur
la page de vérification.
\end{attention}

% ##########################################################
\part{Portail d'administration}
% ##########################################################

\section{Accès et rubriques}
L'accès est réservé aux comptes \etat{ADMIN} ; toute page \chemin{/admin} non authentifiée
redirige vers la connexion. Le menu latéral est organisé en groupes :
\begin{itemize}[leftmargin=1.4em]
  \item \textbf{Pilotage} : Tableau de bord, Réservations, Paiements, Utilisateurs, E-mails envoyés.
  \item \textbf{Offres} : Vols, Hébergements, Biens (chambres), Transports, Locations de véhicules, Activités.
  \item \textbf{Référentiels} : Compagnies aériennes, Prestataires transport, Marques, Modèles, Catégories de véhicules.
  \item \textbf{Configuration} : paramètres de l'agence.
  \item \textbf{Contenu} : Pages d'information, FAQ, Messages de contact.
\end{itemize}

\section{Tableau de bord}
Écran d'accueil agrégeant en un coup d'\oe{}il : le \textbf{chiffre d'affaires} consolidé, la
répartition des \textbf{réservations par statut} et \textbf{par service}, l'évolution récente du
volume, le nombre de \textbf{clients} et le nombre d'offres par service.

\section{Gestion des réservations}

\subsection{Consulter et rechercher}
La page \textbf{Réservations} liste tous les dossiers (référence, client, service, montant,
statut) avec trois outils :
\begin{itemize}[leftmargin=1.4em]
  \item \textbf{Recherche} libre : référence, nom du client, e-mail ou intitulé.
  \item \textbf{Filtre statut} ; \textbf{Filtre service}.
\end{itemize}

\subsection{Cycle de vie d'une réservation}
\begin{center}
\begin{longtable}{>{\raggedright}p{4.2cm} p{9cm}}
\toprule
\textbf{Statut} & \textbf{Signification} \\
\midrule
\endhead
\etat{PENDING\_PAYMENT} & Réservation créée, en attente de règlement. \\
\etat{PAID} & Paiement reçu (en ligne ou encaissé manuellement). \\
\etat{AWAITING\_CONFIRMATION} & Paiement validé, confirmation du prestataire en cours. \\
\etat{CONFIRMED} & Dossier confirmé : billet / bon de réservation disponible. \\
\etat{MODIFIED} & Dossier modifié après confirmation. \\
\etat{CANCELLED} & Réservation annulée. \\
\etat{REFUNDED} & Réservation remboursée. \\
\bottomrule
\end{longtable}
\end{center}

\begin{attention}
Seules les transitions \textbf{autorisées} depuis le statut courant sont proposées dans le
sélecteur, ce qui évite les manipulations incohérentes. Le passage en \etat{CANCELLED} ou
\etat{REFUNDED} remet les disponibilités en stock le cas échéant.
\end{attention}

\section{Détail d'une réservation (vue 360\textdegree)}
En cliquant sur une référence, on ouvre la fiche complète :
\begin{itemize}[leftmargin=1.4em]
  \item \textbf{Informations client} : nom, e-mail, coordonnées.
  \item \textbf{Facturation} : service, quantité, devise, montant total.
  \item \textbf{Voyageurs / participants.}
  \item \textbf{Historique des paiements} du dossier.
  \item \textbf{QR code} du billet (dès que le statut le permet).
  \item \textbf{Documents PDF} : confirmation, récapitulatif, facture, bon de réservation.
  \item \textbf{Actions} : encaissement manuel et renvoi d'e-mail.
\end{itemize}

\subsection{Encaissement manuel}
Pour les règlements reçus hors ligne (Orange Money, MTN MoMo, WhatsApp, espèces, virement\dots) :
\begin{enumerate}[leftmargin=1.6em]
  \item Ouvrir un dossier au statut \etat{PENDING\_PAYMENT}.
  \item Dans le bloc \textbf{Encaissement}, choisir le moyen reçu et valider.
  \item Un paiement \etat{SUCCEEDED} (fournisseur « MANUAL ») est créé et la réservation passe
        automatiquement en \etat{PAID}.
\end{enumerate}
\begin{astuce}
Le bloc d'encaissement n'apparaît que pour les dossiers \textit{en attente de paiement}.
\end{astuce}

\subsection{Renvoi d'e-mail}
Ré-expédition d'une communication au client, avec les \textbf{pièces jointes PDF} le cas échéant :
\begin{center}
\begin{longtable}{>{\raggedright}p{5cm} p{8.2cm}}
\toprule
\textbf{Type d'e-mail} & \textbf{Contenu} \\
\midrule
\endhead
En attente de paiement & Rappel de règlement. \\
Paiement reçu & Accusé de réception + facture. \\
Confirmé & Confirmation + billet / bon de réservation. \\
Modifié & Notification de modification. \\
Annulé & Avis d'annulation. \\
Remboursé & Avis de remboursement. \\
\bottomrule
\end{longtable}
\end{center}

\section{Créer une réservation pour un client}
Vente au comptoir ou par téléphone, sans passage du client par le site :
\begin{enumerate}[leftmargin=1.6em]
  \item Page Réservations $\rightarrow$ \textbf{Nouvelle réservation}.
  \item Sélectionner le \textbf{client}, le \textbf{service}, puis l'\textbf{offre publiée}.
  \item Pour un hébergement ou une location : renseigner les \textbf{dates} et la quantité.
  \item Saisir les \textbf{voyageurs / participants} (un nom par ligne).
  \item Valider : la réservation est créée en \etat{PENDING\_PAYMENT}, un e-mail est envoyé au
        client, et elle devient encaissable depuis son détail.
\end{enumerate}
\begin{attention}
Seules les \textbf{offres publiées} (et, pour vols/transports/activités, à venir) sont
proposées. Si aucune offre n'existe pour un service, la liste l'indique.
\end{attention}

\section{Paiements}
La page \textbf{Paiements} centralise tous les encaissements (en ligne et manuels) avec le
rapprochement client $\rightarrow$ réservation : indicateurs de synthèse, recherche, filtres par
\textbf{statut} et par \textbf{moyen}, et \textbf{export CSV}.

\section{Utilisateurs et rôles}
Liste des comptes (nom, e-mail, rôle, langue, nombre de réservations, date d'inscription), avec
recherche et export CSV. Le rôle (\etat{CLIENT} ou \etat{ADMIN}) se change depuis la liste.
\begin{attention}
\textbf{Anti-verrouillage :} un administrateur ne peut pas rétrograder son propre compte : il
reste toujours au moins un administrateur actif.
\end{attention}

\section{Catalogue des offres et référentiels}
La gestion du catalogue suit un modèle unifié. Chaque entité se liste, se crée, se modifie et
(selon le cas) se publie depuis \chemin{/admin/<entité>}.

\subsection{Offres (publiables)}
Vols, Hébergements, Biens (chambres/suites), Transports, Locations de véhicules, Activités.
Chaque offre a un statut \etat{DRAFT} (brouillon) ou \etat{PUBLISHED} (publié). Seules les offres
\textbf{publiées} sont visibles côté client et réservables.

\subsection{Référentiels}
Compagnies aériennes, prestataires de transport, marques / modèles / catégories de véhicules.
Ils alimentent les listes déroulantes lors de la création des offres (ex. la compagnie d'un vol).

\section{Contenu éditorial}
\begin{itemize}[leftmargin=1.4em]
  \item \textbf{Pages d'information} : contenu des pages du site (ex. infos pratiques).
  \item \textbf{FAQ} : questions/réponses affichées publiquement.
  \item \textbf{Messages de contact} : demandes reçues via le formulaire du site.
\end{itemize}

\section{E-mails envoyés et configuration}
\begin{itemize}[leftmargin=1.4em]
  \item \textbf{E-mails envoyés} : journal des communications adressées aux clients.
  \item \textbf{Configuration} : paramètres de l'agence (coordonnées, horaires, réseaux, etc.).
\end{itemize}

\section{Exports CSV (comptabilité)}
Trois exports via le bouton \textbf{Exporter CSV} : \textbf{Réservations}, \textbf{Paiements},
\textbf{Utilisateurs}. Fichiers encodés en \textbf{UTF-8 avec BOM} et séparés par
\textbf{point-virgule} : ouverture directe dans Excel (configuration française) sans souci
d'accents. Accès réservé aux administrateurs.

\vspace{1.2cm}
\begin{center}
\color{taStone}\small Fin du guide
\end{center}

\end{document}
